%=======================================================================
%  CSE558 Data Science - Midterm Project Report Template
%  Monsoon 2026, IIIT-Delhi
%
%  How to use this template
%   1. Set your project title in the EDIT THIS block below.
%   2. Replace everything in LIGHT BLUE with your own text and numbers.
%      The numbers in the example are made up for illustration.
%   3. Grey italic lines are instructions. Delete them as you write, or
%      hide them all at once with \showhintsfalse.
%   4. Upload your plots to the project and put their file names in
%      \figfile{...}{...}. Until a file exists, a labelled box is shown.
%
%  Sections 1-5 follow items 1-5 of the report guidelines; Section 6 is
%  the limitations of your analysis. The link to your code (item 6) goes
%  at the top. At most 8 pages.
%  Compile with pdfLaTeX (Overleaf's default setting works), twice.
%=======================================================================
\documentclass[10pt,a4paper,twocolumn]{article}

%========================= EDIT THIS =========================
% Your project title: it sets the heading AND the PDF's title.
\newcommand{\projecttitle}{What separates a fraudulent card transaction from a normal one?
  Evidence from 1.85 million simulated transactions}
\newif\ifshowhints
\showhintsfalse     % \showhintsfalse hides every grey instruction line
%=============================================================

%-------------------------------------------- fonts (no need to edit)
\usepackage{iftex}
\ifPDFTeX
  \usepackage[T1]{fontenc}
  \usepackage[utf8]{inputenc}
  \IfFileExists{carlito.sty}%
    {\usepackage[sfdefault]{carlito}}%              Calibri look
    {\usepackage[scaled=0.92]{helvet}\renewcommand{\familydefault}{\sfdefault}}%
  \usepackage{courier}
\else
  \usepackage[no-math]{fontspec}
  \setmainfont{Carlito}[Extension=.ttf, UprightFont=*-Regular, BoldFont=*-Bold,
    ItalicFont=*-Italic, BoldItalicFont=*-BoldItalic]
\fi

%----------------------------------------------------------- packages
\usepackage[margin=18mm,bottom=20mm,columnsep=7mm]{geometry}
\usepackage{amsmath}
\usepackage{xcolor}
\usepackage{graphicx}
\usepackage{tabularx}
\usepackage{array}
\usepackage{booktabs}
\usepackage{enumitem}
\usepackage{titlesec}
\usepackage{needspace}
\usepackage[font=small,labelfont=bf]{caption}
\usepackage[hidelinks]{hyperref}
\hypersetup{pdftitle={\projecttitle},
            pdfsubject={CSE558 Data Science, Monsoon 2026: Midterm Project Report}}

%------------------------------------------------------------- layout
\definecolor{eg}{HTML}{7396C4}         % worked example (light blue)
\definecolor{hintgrey}{gray}{0.58}     % instructions (light grey)

\titleformat{\section}{\normalfont\large\bfseries}{\thesection.}{0.5em}{}
\titlespacing*{\section}{0pt}{12pt}{3pt}
\titleformat{\subsection}{\normalfont\normalsize\bfseries}{\thesubsection}{0.5em}{}
\titlespacing*{\subsection}{0pt}{10pt}{2pt}
\newcommand{\sectionbreak}{\needspace{10\baselineskip}}
\newcommand{\subsectionbreak}{\needspace{10\baselineskip}}

\setlength{\parindent}{0pt}
\setlength{\parskip}{4pt}
\setlist{leftmargin=1.2em, itemsep=2pt, topsep=2pt, parsep=0pt}
\renewcommand{\arraystretch}{1.25}
\setlength{\tabcolsep}{4pt}
\newcolumntype{L}[1]{>{\raggedright\arraybackslash}p{#1}}
\newcolumntype{Y}{>{\raggedright\arraybackslash}X}
\newcolumntype{K}[1]{>{\raggedright\arraybackslash\bfseries\boldmath}p{#1}}  % bold labels

%--------------------------------------------------------- shorthands
% Instruction line. Delete these as you fill the template in.
\newcommand{\hint}[1]{\ifshowhints{\footnotesize\itshape\color{hintgrey}#1\par}\fi}
% Worked example: inline or in a table cell. REPLACE the contents.
\newcommand{\eg}[1]{\textcolor{eg}{#1}}
% Worked example running over a whole paragraph.
\newcommand{\egblock}[1]{{\color{eg}#1\par}}
% Bold label followed by an example.
\newcommand{\field}[2]{\textbf{#1}~\eg{#2}\par}
% Part a)-c) of a test, and a small bold heading.
\newcommand{\tpart}[2]{\par\needspace{12\baselineskip}\smallskip\textbf{#1)~#2}\par}
\newcommand{\minihead}[1]{\par\needspace{8\baselineskip}\smallskip\textbf{#1}\par}
% Column and file names.
\newcommand{\code}[1]{\texttt{#1}}
% \figfile{file}{height}: your plot, or a labelled box until the file exists.
\newcommand{\figfile}[2]{%
  \IfFileExists{#1}%
    {\includegraphics[width=\linewidth,height=#2,keepaspectratio]{#1}}%
    {\fbox{\parbox[c][#2][c]{\dimexpr\linewidth-2\fboxsep-2\fboxrule\relax}%
      {\centering\small\itshape\color{hintgrey}Your plot: \detokenize{#1}}}}}

%=======================================================================
\begin{document}

\twocolumn[%
{\small\bfseries CSE558 Data Science\ \ \textbullet\ \ Midterm Project Report\ \ \textbullet\ \ Monsoon 2026\par}
\vspace{4pt}
{\LARGE\bfseries\raggedright\projecttitle\par}
\vspace{6pt}
\rule{\linewidth}{1pt}

{\small
\begin{tabularx}{\linewidth}{@{}K{2.4cm} Y@{}}
\toprule
Group name & Null \& Void \\
Dataset & kaggle.com/datasets/kartik2112/fraud-detection \\
Code & \href{https://github.com/DarshanBothra/cc-fraud-detection/tree/main/midterm}{github.com/DarshanBothra/cc-fraud-detection/tree/main/midterm} (code in \code{midterm/src}, README in \code{midterm/README.md}) \\
\bottomrule
\end{tabularx}}

\vspace{2pt}
{\small
\begin{tabularx}{\linewidth}{@{}L{0.5cm} L{3.4cm} L{2.3cm} L{4.8cm} Y@{}}
\toprule
\textbf{\#} & \textbf{Name} & \textbf{Roll no.} & \textbf{Institute email} & \textbf{Contribution} \\
\midrule
1 & Darshan Bothra & 2024176 & darshan24176@iiitd.ac.in & Dataset choice, repository, classifier design (1, 5) \\
2 & Piyush Keshan & 2023376 & piyush23376@iiitd.ac.in & Test plan, code, hypothesis tests, report (2--6) \\
3 & Virat Yadav & 2024622 & virat24622@iiitd.ac.in & Group and topic registration, review (1, 6) \\
\bottomrule
\end{tabularx}}
\vspace{12pt}
]

%-----------------------------------------------------------------------
\section{Dataset}

\textbf{Source:} Kaggle, ``Credit Card Transactions Fraud Detection Dataset'' (Kartik Shenoy, CC0),
generated with Brandon Harris's Sparkov simulator. One row is one card transaction by one of 999
customers at one of 693 merchants, 1~Jan~2019 to 31~Dec~2020. It comes as \code{fraudTrain.csv}
(1{,}296{,}675 rows) and \code{fraudTest.csv} (555{,}719 rows; later months), which we combine.
The response \code{is\_fraud} marks 9{,}651 frauds (0.52\%), so the classes are heavily imbalanced.

\textbf{What is raw about it:} the timestamp and date of birth are text; customer and merchant
locations are four separate coordinates rather than a distance; amounts are extremely skewed
(median \$47, maximum \$28{,}949); and the file carries names, street addresses and card numbers that
identify people.

{\small
\begin{tabularx}{\linewidth}{@{}L{3.0cm} Y Y@{}}
\toprule
 & \textbf{Raw file} & \textbf{After encoding} \\
\midrule
Instances (rows) & 1{,}852{,}394 & 1{,}852{,}394 \\
Attributes (columns) & 23 & 28 \\
Instances $\times$ attributes & 42{,}605{,}062 & 51{,}867{,}032 \\
\bottomrule
\end{tabularx}}

\textbf{Size check:} $1{,}852{,}394\times28=51{,}867{,}032\ge3{,}000{,}000$, so the requirement is
met. The 28 columns are 13 numeric or derived features, the response, and 14 one-hot category
columns.

\minihead{Key features}
{\footnotesize
\begin{tabularx}{\linewidth}{@{}L{2.2cm} L{1.5cm} Y@{}}
\toprule
\textbf{Feature} & \textbf{Type} & \textbf{Real-world meaning} \\
\midrule
\code{is\_fraud} & nominal & transaction was fraudulent (1) or not (0); our response \\
\code{amt} & continuous & charge in US dollars \\
\code{category} & nominal & merchant type, 14 levels (e.g.\ \code{grocery\_pos} = in-store groceries, \code{shopping\_net} = online shopping) \\
\code{hour} & discrete & hour of the transaction, from the timestamp \\
\code{age} & continuous & cardholder's age at the time of purchase, from date of birth \\
\code{distance\_km} & continuous & distance from cardholder's home to the merchant (derived) \\
\code{city\_pop} & discrete & population of the cardholder's city \\
\code{gender} & nominal & cardholder's gender \\
\bottomrule
\end{tabularx}}

%-----------------------------------------------------------------------
\section{Data preparation}

{\footnotesize
\begin{tabularx}{\linewidth}{@{}L{1.9cm} L{2.3cm} Y@{}}
\toprule
\textbf{Issue found} & \textbf{What we did} & \textbf{Why, and not the alternative} \\
\midrule
Missing values, duplicates
 & checked all 23 columns and the transaction id: 0 missing, 0 duplicates
 & nothing to impute or drop, so no rows are lost and no values are invented \\
Timestamp stored as text
 & parsed; derived hour, night flag (22:00--03:59), day of week, month
 & time of day is one of our questions; the raw string cannot be compared \\
Date of birth, not age
 & age at transaction $=(\text{date}-\text{dob})/365.25$
 & age at purchase is what matters; age today would shift every row by up to two years \\
Four coordinates, no distance
 & haversine distance between home and merchant, in km
 & our Test 5 is about distance; no single raw column holds it \\
\code{amt} skewness 40.8; 95{,}054 points above the IQR fence
 & kept every row; analysed $\log(1+\code{amt})$
 & the large amounts are where fraud is (median fraud \$390 vs \$47), so deleting or capping them would remove frauds \\
\code{category}, 14 levels
 & one-hot encoded
 & nominal, no order; needed for the classifier and the size check \\
Names, street, card number, transaction id, merchant, job, city, zip, state
 & dropped
 & they identify people or are high-cardinality identifiers with no role in our questions \\
\bottomrule
\end{tabularx}}

%-----------------------------------------------------------------------
\section{What changed after preparation}

{\footnotesize
\begin{tabularx}{\linewidth}{@{}L{2.9cm} L{1.5cm} Y@{}}
\toprule
\textbf{Measure} & \textbf{Before} & \textbf{After} \\
\midrule
Rows & 1{,}852{,}394 & 1{,}852{,}394 (none lost) \\
Attributes & 23 & 28 (identifiers and index dropped; 5 derived; category one-hot) \\
Missing values, duplicates & 0, 0 & 0, 0 \\
Skewness of amount & 40.81 & $-0.30$ (log scale) \\
Points beyond the IQR fence & 95{,}054 & 1{,}220 (log scale) \\
Age, distance ranges & not available & 13.9--96.2 years; 0.02--152.1 km \\
\bottomrule
\end{tabularx}}

The log transform turns a long right tail into a near-symmetric distribution
(Figure~\ref{fig:prep}), which is what makes a $t$-test on amounts reasonable. Deriving age, hour and
distance turned three text or coordinate fields into the quantities our questions are about. All
derived values are in plausible ranges, so no further cleaning was needed.

\begin{figure}[!ht]
  \figfile{../figures/prep_before_after.pdf}{2.6cm}
  \caption{Amount before (left, clipped at \$1{,}500 for display) and after (right) the log
  transform.}
  \label{fig:prep}
\end{figure}

%-----------------------------------------------------------------------
\section{Basic questions and problem statement}

\begin{figure}[!ht]
  \figfile{../figures/eda_log_amt.pdf}{3.4cm}
  \caption{Frauds are far larger: median \$390 against \$47 for legitimate transactions.}
  \label{fig:amt}
\end{figure}

\begin{figure}[!ht]
  \figfile{../figures/eda_hour.pdf}{3.4cm}
  \caption{The fraud rate is about 2.6\% at 22--23h and 1.3\% at 0--3h, against about 0.1\% in
  every other hour (night hours in red).}
  \label{fig:hour}
\end{figure}

\begin{figure}[!ht]
  \figfile{../figures/eda_category.pdf}{4.0cm}
  \caption{Online shopping, misc.\ online and in-store groceries have fraud rates of 1.3--1.6\%,
  about 2.4--3$\times$ the overall 0.52\% (dashed).}
  \label{fig:cat}
\end{figure}

\begin{figure}[!ht]
  \figfile{../figures/eda_age.pdf}{3.2cm}
  \caption{Fraud victims are slightly older (mean 48.8 against 46.3 years), and the shapes differ.}
  \label{fig:age}
\end{figure}

\begin{figure}[!ht]
  \figfile{../figures/eda_distance.pdf}{3.2cm}
  \caption{Home-to-merchant distance looks the same for both classes (mean 76.3 against
  76.1~km).}
  \label{fig:dist}
\end{figure}

\minihead{Basic questions}
\begin{enumerate}[label=\arabic*.]
  \item How common is fraud, and is it spread evenly over the two years?
  \item Which merchant categories see the most fraud?
  \item Are fraudulent transactions more expensive?
\end{enumerate}

\minihead{Our problem statements}
\noindent\fbox{\begin{minipage}{\dimexpr\linewidth-2\fboxsep-2\fboxrule\relax}
\textbf{P1 (primary).} Do fraudulent transactions differ from legitimate ones in \emph{how much},
\emph{when} and \emph{where} money is spent (amount, time of day, merchant category) by margins large
enough to matter to a bank's fraud team?

\smallskip
\textbf{P2.} Are the cardholders hit by fraud different: are they older, and does fraud happen
farther from their home?
\end{minipage}}

\textbf{How they are novel:} the basic questions describe fraud rates. Ours fix, before testing, the
smallest difference that would matter, test on a power-calculated sample rather than on all 1.85
million rows, use an engineered distance feature, and then check (Section~5.6) whether each tested
signal is one a classifier actually needs.

%-----------------------------------------------------------------------
\section{Statistical inferences}

\textbf{Rows tested.} Not the full data: on 1.85 million rows every test would reject $H_0$ for
differences far too small to matter. We draw a \textbf{stratified random sample of 2{,}000 fraud and
2{,}000 legitimate transactions} (seed 42, without replacement). A simple random sample would not
work: 4{,}000 random rows hold only about 21 frauds. The size comes from a \textbf{power
calculation}: a separate pilot sample of 500 per class gave the standard deviations, and
$n=2(z_{1-\alpha}+z_{1-\beta})^2\sigma^2/\Delta^2$ per group at $\alpha=0.01$ (Bonferroni for five
tests) and power $0.8$ requires 263 (amount), 383 (night share), 810 (age) and 786 (distance) per
group, so 2{,}000 covers every test. The sample is disjoint from the pilot and representative: within
each class, mean amount, age, distance and night share are within 6\% of the full data, and no
category share differs by more than 1.1 points. Sample ids are saved in the repository.

\textbf{Consequence.} Fixing 2{,}000 per class makes the sample's fraud rate (50\%) artificial, so
Tests 2 and 3 compare how night-time and categories are distributed \emph{within} each class.
Independence is symmetric, so this answers the same question, and odds ratios stay valid (the
case--control argument). Each test uses $\alpha=0.05$; we also report whether it survives
Bonferroni's $0.01$.

%------------------------------------------------------------ Test 1
\subsection*{Test 1 (P1): are fraudulent transactions larger?}

\tpart{a}{Hypotheses, written before running the test}
{\small
\begin{tabularx}{\linewidth}{@{}K{3.4cm} Y@{}}
\toprule
$H_0$ & $\mu_F-\mu_L\le0$ \\
$H_1$ & $\mu_F-\mu_L>0$ \\
Significance level & $\alpha=0.05$ (also 0.01) \\
Side & one-sided (right) \\
Test & Welch $t$-test on $\log(1+\code{amt})$ \\
Smallest difference that matters & fraud amounts $1.5\times$ larger ($\Delta=\ln1.5=0.405$) \\
\bottomrule
\end{tabularx}}

$\mu_F$, $\mu_L$: mean log amount of fraudulent and legitimate transactions. One-sided, because P1
expects frauds to be larger. On the log scale a difference is a ratio of typical amounts; a fraud
team would not change a rule for less than a 50\% gap.

\tpart{b}{Other tests we considered}
Two independent groups, a continuous response and unknown $\sigma$ call for a two-sample $t$-test.
An $F$-test rejects equal variances ($F=1.66$, df $(1999,1999)$, $p<10^{-300}$), so we use Welch's
version.
{\footnotesize
\begin{tabularx}{\linewidth}{@{}L{2.4cm} Y@{}}
\toprule
\textbf{Also considered} & \textbf{Why we did not use it} \\
\midrule
Two-sample $z$-test & needs the population $\sigma$; at $n=2{,}000$ it gives the same answer \\
Pooled $t$-test & assumes equal variances, which the $F$-test rejects. Run anyway: same $t$ and decision \\
Mann--Whitney $U$ & no normality assumption, but tests ranks, not the mean ratio we want; run as a check, it agrees ($p\approx10^{-285}$) \\
\bottomrule
\end{tabularx}}

\tpart{c}{Result}
{\small
\begin{tabularx}{\linewidth}{@{}K{3.4cm} Y@{}}
\toprule
Estimate & $+2.02$ on the log scale: fraud amounts $7.5\times$ larger (means 5.56 vs 3.54) \\
95\% confidence interval & $1.93$ to $2.11$, i.e.\ $6.9\times$ to $8.2\times$ \\
Test statistic (d.f.) & $t=43.67$ (3{,}763.8) \\
p-value & $<10^{-300}$ \\
Decision & reject $H_0$ (also at 0.01) \\
\bottomrule
\end{tabularx}}

Frauds are about 7.5 times larger than legitimate transactions. The whole interval is far above the
$1.5\times$ that matters, so the difference is real and large.

%------------------------------------------------------------ Test 2
\subsection*{Test 2 (P1): is fraud concentrated at night?}

\tpart{a}{Hypotheses, written before running the test}
{\small
\begin{tabularx}{\linewidth}{@{}K{3.4cm} Y@{}}
\toprule
$H_0$ & $p_F-p_L\le0$ \\
$H_1$ & $p_F-p_L>0$ \\
Significance level & $\alpha=0.05$ (also 0.01) \\
Side & one-sided (right) \\
Test & two-proportion $z$-test \\
Smallest difference that matters & 10 percentage points \\
\bottomrule
\end{tabularx}}

$p_F$, $p_L$: share of fraudulent and of legitimate transactions made at night (22:00--03:59).
One-sided, because P1 expects fraud to favour the night. We judged a smaller shift too weak to base
a time-of-day rule on.

\tpart{b}{Other tests we considered}
Two independent groups and a binary outcome (night or not) call for a two-proportion $z$-test; with
2{,}000 per group the normal approximation is safe.
{\footnotesize
\begin{tabularx}{\linewidth}{@{}L{2.4cm} Y@{}}
\toprule
\textbf{Also considered} & \textbf{Why we did not use it} \\
\midrule
$\chi^2$ test on the $2\times2$ table & equivalent ($\chi^2=1544.6=z^2$) but two-sided only \\
Fisher's exact test & exact, but unnecessary with 2{,}000 per group; it agrees \\
Risk ratio & cannot be estimated from a class-stratified sample; we report the odds ratio instead \\
\bottomrule
\end{tabularx}}

\tpart{c}{Result}
{\small
\begin{tabularx}{\linewidth}{@{}K{3.4cm} Y@{}}
\toprule
Estimate & $+61.9$ points (85.4\% of frauds vs 23.5\% of legitimate at night); odds ratio 19.0 \\
95\% confidence interval & $+59.5$ to $+64.3$ points; odds ratio 16.2 to 22.4 \\
Test statistic & $z=39.30$ \\
p-value & $<10^{-300}$ \\
Decision & reject $H_0$ (also at 0.01) \\
\bottomrule
\end{tabularx}}

Most frauds happen at night, against under a quarter of legitimate spending. The interval is six
times the 10 points that matter.

%------------------------------------------------------------ Test 3
\subsection*{Test 3 (P1): does fraud depend on the merchant category?}

\tpart{a}{Hypotheses, written before running the test}
{\small
\begin{tabularx}{\linewidth}{@{}K{3.4cm} Y@{}}
\toprule
$H_0$ & category is independent of \code{is\_fraud} (same category mix in both classes) \\
$H_1$ & the category mix differs between the classes \\
Significance level & $\alpha=0.05$ (also 0.01) \\
Side & not applicable ($\chi^2$, upper tail) \\
Test & $\chi^2$ test of independence, $14\times2$ table, df $=13$ \\
Smallest difference that matters & a category with odds ratio $\ge2$ (fraud vs legitimate) \\
\bottomrule
\end{tabularx}}

An odds ratio of 2 means a category is twice as over-represented among frauds; below that, flagging
a whole category would not be worth the false alarms.

\tpart{b}{Other tests we considered}
Two nominal variables call for a $\chi^2$ test of independence; the smallest expected count is 49.5,
so the approximation is valid.
{\footnotesize
\begin{tabularx}{\linewidth}{@{}L{2.4cm} Y@{}}
\toprule
\textbf{Also considered} & \textbf{Why we did not use it} \\
\midrule
G-test (likelihood ratio) & asymptotically equivalent; run as a check: $G=778.7$, same decision \\
Fisher's exact test & impractical for a $14\times2$ table with 4{,}000 rows, and unnecessary with large expected counts \\
14 separate proportion tests & multiple comparisons; we use per-category odds ratios with intervals only as a follow-up \\
\bottomrule
\end{tabularx}}

\tpart{c}{Result}
{\small
\begin{tabularx}{\linewidth}{@{}K{3.4cm} Y@{}}
\toprule
Estimate & Cramér's $V=0.43$ (strong). Odds ratios: \code{shopping\_net} 3.71, \code{grocery\_pos} 3.22, \code{misc\_net} 2.48; lowest \code{home} 0.18 \\
95\% confidence interval & 3.05--4.52; 2.67--3.88; 1.96--3.14; 0.13--0.26 \\
Test statistic (d.f.) & $\chi^2=742.0$ (13) \\
p-value & $3.6\times10^{-150}$ \\
Decision & reject $H_0$ (also at 0.01) \\
\bottomrule
\end{tabularx}}

Fraud depends strongly on category. Online shopping and in-store groceries clear the odds-ratio-2 bar
with their whole interval; misc.\ online very nearly does.

%------------------------------------------------------------ Test 4
\subsection*{Test 4 (P2): are fraud victims older?}

\tpart{a}{Hypotheses, written before running the test}
{\small
\begin{tabularx}{\linewidth}{@{}K{3.4cm} Y@{}}
\toprule
$H_0$ & $\mu_F-\mu_L=0$ (mean age) \\
$H_1$ & $\mu_F-\mu_L\ne0$ \\
Significance level & $\alpha=0.05$ (also 0.01) \\
Side & two-sided \\
Test & Welch $t$-test \\
Smallest difference that matters & 3 years \\
\bottomrule
\end{tabularx}}

Two-sided, because we had no prior direction. We judged a gap under 3 years too small to justify
age-specific fraud rules.

\tpart{b}{Other tests we considered}
Two independent groups, a continuous response, unknown $\sigma$: a two-sample $t$-test, Welch's
version fixed in advance (SDs 18.7 and 18.1).
{\footnotesize
\begin{tabularx}{\linewidth}{@{}L{2.4cm} Y@{}}
\toprule
\textbf{Also considered} & \textbf{Why we did not use it} \\
\midrule
Mann--Whitney $U$ & tests ranks rather than mean age; agrees ($p=4\times10^{-6}$) \\
Kolmogorov--Smirnov & tests any difference in shape, not the mean; run as a complement: $D=0.107$, $p=2\times10^{-10}$ \\
TOST equivalence test & used as a follow-up, to ask whether the difference is below 3 years \\
\bottomrule
\end{tabularx}}

\tpart{c}{Result}
{\small
\begin{tabularx}{\linewidth}{@{}K{3.4cm} Y@{}}
\toprule
Estimate & $+2.6$ years (49.2 vs 46.7 in the sample) \\
95\% confidence interval & $+1.4$ to $+3.7$ years \\
Test statistic (d.f.) & $t=4.40$ (3{,}994.2) \\
p-value & $1.1\times10^{-5}$ \\
Decision & reject $H_0$ (also at 0.01) \\
\bottomrule
\end{tabularx}}

Fraud victims are older, but the interval contains the 3 years that would matter, and an equivalence
test cannot rule out a 3-year gap either (TOST $p=0.23$). The difference is real; whether it matters
is unresolved. The KS test adds that the age \emph{distributions} differ in shape (Figure~\ref{fig:age}).

%------------------------------------------------------------ Test 5
\subsection*{Test 5 (P2): does fraud happen farther from home?}

\tpart{a}{Hypotheses, written before running the test}
{\small
\begin{tabularx}{\linewidth}{@{}K{3.4cm} Y@{}}
\toprule
$H_0$ & $\mu_F-\mu_L=0$ (mean distance) \\
$H_1$ & $\mu_F-\mu_L\ne0$ \\
Significance level & $\alpha=0.05$ (also 0.01) \\
Side & two-sided \\
Test & Welch $t$-test on \code{distance\_km} \\
Smallest difference that matters & 5 km \\
\bottomrule
\end{tabularx}}

Two-sided: fraud could be farther (stolen card used elsewhere) or nearer. Under 5~km is within one
town, too small for a location rule.

\tpart{b}{Other tests we considered}
Again two independent groups and a continuous response with unknown $\sigma$, so Welch's $t$-test
(SDs 28.8 and 28.7).
{\footnotesize
\begin{tabularx}{\linewidth}{@{}L{2.4cm} Y@{}}
\toprule
\textbf{Also considered} & \textbf{Why we did not use it} \\
\midrule
Mann--Whitney $U$ & tests ranks; borderline ($p=0.046$), not significant at 0.01 \\
Kolmogorov--Smirnov & tests shape; $p=0.064$ \\
TOST equivalence test & used as a follow-up, because failing to reject alone would prove nothing \\
\bottomrule
\end{tabularx}}

\tpart{c}{Result}
{\small
\begin{tabularx}{\linewidth}{@{}K{3.4cm} Y@{}}
\toprule
Estimate & $+1.8$ km (76.6 vs 74.8 km) \\
95\% confidence interval & $-0.02$ to $+3.54$ km \\
Test statistic (d.f.) & $t=1.94$ (3{,}998.0) \\
p-value & 0.053 \\
Decision & fail to reject $H_0$ \\
\bottomrule
\end{tabularx}}

No evidence that fraud happens farther from home. Failing to reject is not proof of equality, so we
ran TOST with bounds $\pm5$~km: it rejects non-equivalence ($p=0.0002$). The sample was powered for
5~km, so any difference is practically negligible.

%------------------------------------------------------------ summary
\subsection*{5.6 Robustness, and do the signals matter in a classifier?}

All five decisions at $\alpha=0.01$ are the same on five further independent samples (seeds
1001--1005). As an extra, a $\chi^2$ goodness-of-fit test shows fraud is not spread over the days
of the week like legitimate spending ($\chi^2=88.5$, df 6, $p=6\times10^{-17}$).

Since the data is labelled, we fit binary classifiers on \code{fraudTrain.csv}
(threshold and settings chosen on a 20\% validation split) and evaluate once on
\code{fraudTest.csv} (555{,}719 rows, 2{,}145 frauds). We report precision, recall, F1 and PR-AUC,
not accuracy: calling everything legitimate already scores 99.6\%.

{\footnotesize
\begin{tabularx}{\linewidth}{@{}Y rrrr@{}}
\toprule
\textbf{Model (test file)} & \textbf{Prec.} & \textbf{Rec.} & \textbf{F1} & \textbf{PR-AUC} \\
\midrule
Logistic regression, linear & 0.32 & 0.58 & 0.41 & 0.39 \\
Logistic regression, splines & 0.64 & 0.70 & 0.67 & 0.70 \\
Decision tree (depth 12) & 0.87 & 0.78 & 0.82 & 0.85 \\
\bottomrule
\end{tabularx}}

\textbf{McNemar's test} ($H_0$: equal error rates): the tree is right where spline logistic
regression is wrong 944 times, the reverse 180 times; $\chi^2=517.9$, $p\approx10^{-114}$. A paired
bootstrap puts the tree's F1 advantage at 0.14--0.17.

\textbf{Ablation:} dropping each tested feature group and refitting, the fall in test PR-AUC (tree;
spline logistic regression in brackets) is: amount 0.61 (0.48), category 0.21 (0.32), time of day
0.11 (0.11), age 0.04 (0.01), distance 0.00 (0.00). The features with large, practically important
differences in Tests 1--3 are the ones the models cannot do without; age, significant but small,
adds little; distance, shown negligible by TOST, adds nothing.

{\footnotesize
\begin{tabularx}{\linewidth}{@{}L{1.3cm} L{2.6cm} Y L{1.0cm}@{}}
\toprule
\textbf{Test} & \textbf{Estimate (95\% CI)} & \textbf{$p$} & \textbf{Matters?} \\
\midrule
1 Amount & $7.5\times$ ($6.9$--$8.2$) & $<10^{-300}$ & yes \\
2 Night & $+61.9$ pts ($59.5$--$64.3$) & $<10^{-300}$ & yes \\
3 Category & $V=0.43$; OR up to 3.7 & $10^{-150}$ & yes \\
4 Age & $+2.6$ yr ($1.4$--$3.7$) & $10^{-5}$ & unclear \\
5 Distance & $+1.8$ km ($-0.02$--$3.5$) & 0.053 & no \\
\bottomrule
\end{tabularx}}

%-----------------------------------------------------------------------
\section{Limitations}

\begin{itemize}
  \item \textbf{Simulated data.} The transactions come from the Sparkov generator, so our effects
  describe its fraud profiles, not a real bank's.
  \item \textbf{Association, not cause.} Night-time or online purchases are not what makes a
  transaction fraudulent; they are where fraud shows up.
  \item \textbf{Transactions are not fully independent.} 999 customers each make many transactions,
  so the effective sample is smaller than 4{,}000 and our p-values are somewhat optimistic.
  \item \textbf{One test sample.} Results come from one stratified sample, though five further
  samples gave the same decisions.
  \item \textbf{No absolute fraud rates from the sample.} Stratified sampling cannot estimate
  $P(\text{fraud}\mid\cdot)$; such rates are reported only descriptively from the full data.
  \item \textbf{Multiple tests.} Five tests are run; we report both $\alpha=0.05$ and Bonferroni's
  $0.01$, and every decision is the same at both.
\end{itemize}

\end{document}
